
\section{总体研究框架}

本项目围绕"面向量子优化平台的问题规约编译器"这一核心目标，构建包含四个紧密关联的研究方向的系统性框架，如图\ref{fig:framework}所示。

\begin{figure}[h!]
	\centering
	\includegraphics[width=0.85\textwidth]{framework.pdf}
	\captionsetup{justification=centering}
	\caption{总体研究框架：问题规约编译器的理论、方法与系统}
	\label{fig:framework}
\end{figure}

四个研究方向形成从理论到系统的完整链条。\textbf{研究内容一}建立规约编译的理论基础，将规约选择形式化为多项式代价图上的路径优化问题，研究最优路径的判定条件和搜索算法。\textbf{研究内容二}提出测试驱动的规约发现方法，利用大语言模型自动生成候选规约，结合ILP进行正确性验证，实现100+规约规则的规模化发现。\textbf{研究内容三}构建编译器系统，集成规约知识库、路径规划器和后端代码生成器，实现从用户问题到量子硬件的端到端自动映射。\textbf{研究内容四}在代表性量子优化平台上验证编译器生成的规约，评估实际求解效果。

这四个方向形成"理论$\to$方法$\to$系统$\to$验证"的逻辑链条，相互支撑、缺一不可：理论研究（一）为路径选择提供算法基础；方法研究（二）为编译器填充规约知识库；系统研究（三）将理论和方法集成为端到端工具；实验验证（四）在真实量子硬件上检验效果并反馈改进方向。

上述四个方向对应四个可验证的子目标（S1--S4），每个子目标均明确指标、对照基线与验证方案：
\begin{itemize}
\item \textbf{S1（理论）}：建立多变量多项式代价模型与Pareto最优路径搜索算法。\textit{指标}：100+节点规约图的秒级路径搜索。\textit{对照}：穷举搜索与贪心策略。\textit{验证}：Pareto前沿完备性分析与路径质量对比。

\item \textbf{S2（方法）}：发现并验证100+规约规则，覆盖$\geq$15类NP难问题。\textit{指标}：ILP验证通过率与代价多项式精度。\textit{对照}：纯随机搜索与人工设计效率。\textit{验证}：ILP严格证明+热带张量网络\cite{Liu2021PRL}精确求解交叉检验。

\item \textbf{S3（系统）}：构建端到端问题规约编译器。\textit{指标}：千顶点规模映射时间$\leq$分钟量级，量子比特开销。\textit{对照}：手工设计方案与D-Wave Ocean SDK默认嵌入。\textit{验证}：标准基准问题集上的系统对比。

\item \textbf{S4（验证）}：在代表性量子优化平台（里德堡原子阵列或量子退火机）上完成实验验证。\textit{指标}：求解成功率与近似比。\textit{对照}：已有编码方案（如KSG编码\cite{Nguyen2023PRXQuantum}）与平台默认嵌入工具。\textit{验证}：$\geq$5类NP难问题、$\geq$3个规模梯度的系统性实验。
\end{itemize}


\section{研究内容一：规约代价模型与路径优化}

\subsection{研究目标}

将组合优化问题映射到量子硬件通常需要经过多步规约，每步规约产生不同的量子比特开销和能隙特性。本研究的目标是建立规约路径优化的理论基础——给定源问题和目标量子平台，如何在多条可行的规约路径中自动选择资源最优的方案。

\subsection{代价模型}

与经典编译器的常数代价模型不同，规约代价是\textbf{问题规模的函数}。设问题$P$的规模由参数向量$\mathbf{n} = (n_1, n_2, \ldots)$描述，规约$r: P_i \to P_j$的代价可表示为多变量多项式：
\begin{equation}
	\text{cost}_r(\mathbf{n}) = \sum_{k} c_k \cdot n_1^{a_{k1}} n_2^{a_{k2}} \cdots
\end{equation}

例如，将$n$顶点、$m$边的图规约到三角晶格MIS，开销为$O(n^2 + m)$个原子位点；3-SAT到MIS的规约，开销正比于子句数$m$。规约代价直接决定所需的量子比特数量，进而影响实验的可实现性和退相干带来的误差。

\subsection{路径优化问题}

将规约空间建模为有向图$G = (V, E)$，其中顶点$V$为问题类型，边$E$为规约规则。给定源问题$P_s$和目标问题$P_t$，最优路径问题为：
\begin{equation}
	\pi^* = \arg\min_{\pi \in \text{Paths}(P_s, P_t)} \text{TotalCost}(\pi, \mathbf{n})
\end{equation}

\textbf{核心理论问题}：

\textbf{问题1：代价比较。}给定两条路径的代价多项式$f(\mathbf{n})$和$g(\mathbf{n})$，判定是否存在$\mathbf{n}_0$使得对所有$\mathbf{n} > \mathbf{n}_0$有$f(\mathbf{n}) < g(\mathbf{n})$。这涉及多变量多项式的渐进比较理论。

\textbf{问题2：Pareto最优。}当不存在全局最优路径时（即不同规模范围内最优路径不同），需要计算Pareto前沿——在各规模范围内分别最优的路径集合。

\textbf{问题3：高效搜索。}规约图可能包含100+节点和数百条边，需要开发高效的路径枚举和剪枝算法。

\subsection{技术路线}

\textbf{渐进分析框架：}借鉴计算复杂性理论中的渐进分析方法，建立多项式代价的偏序关系。对于单变量情形，利用主项比较；对于多变量情形，研究不同变量增长率假设下的比较策略。

\textbf{Pareto前沿计算：}开发增量式算法维护Pareto最优路径集合。当规约图更新时（新增规约规则），高效更新前沿而无需重新计算。

\textbf{启发式加速：}对于大规模规约图，开发基于问题结构的启发式剪枝策略，在保证解质量的前提下加速搜索。

\subsection{预期成果}

\textbf{理论成果：}建立规约路径优化的形式化框架，给出代价比较的可判定性条件，设计Pareto前沿的增量计算算法。预期发表理论论文1--2篇。

\textbf{算法实现：}开发路径优化算法库，作为编译器的核心模块，支持100+节点规模规约图的秒级路径搜索，开源发布。


\section{研究内容二：LLM驱动的规模化规约发现}

\subsection{研究目标}

构建覆盖主流NP难问题的规约知识库需要发现和验证上百条规约规则。传统人工方法每条规约需要数周至数月的推导与验证，无法满足规模化需求。本研究提出\textbf{LLM驱动的测试驱动发现方法}，将规约发现从"专家手工设计"转变为"人机协作的自动化流水线"。

\subsection{测试驱动的规约发现}

申请人提出规约发现的测试驱动方法，其核心思路是"先严格定义正确性标准，再借助LLM探索解空间"。

第一步是\textbf{规范定义}：形式化定义规约的数学规范，包括源问题$P_s$和目标问题$P_t$的精确定义、规约映射$\phi: \text{Instance}(P_s) \to \text{Instance}(P_t)$（将源问题实例映射到目标哈密顿量）与解提取映射$\psi: \text{Sol}(P_t) \to \text{Sol}(P_s)$（从目标基态恢复源问题最优解）的数学形式，以及正确性性质（解保持性与可提取性）的形式化陈述。

第二步是\textbf{验证体系构建}：在LLM生成规约之前，先建立多层次验证体系——对所有小规模实例（利用热带张量网络\cite{Liu2021PRL}进行精确求解）穷举验证解保持性和可提取性，生成大量随机实例进行统计验证，并针对空图、完全图、稀疏图等极端结构设计边界用例。

第三步是\textbf{LLM生成与迭代}：向大语言模型提供规范定义和验证体系，由LLM生成候选规约方案。验证失败时，将具体的反例（如特定图结构上的解不一致）反馈给LLM进行迭代修正，形成"生成—验证—反馈—改进"的闭环。

第四步是\textbf{严格数学证明}：对通过初步验证的规约，使用ILP求解器进行严格的数学证明，确保正确性源于证明而非仅依赖有限实例的验证。

\subsection{LLM引导策略}

如何有效引导LLM生成高质量规约方案是本研究的关键技术：

\textbf{物理知识编码：}将现有规约设计知识（如KSG编码\cite{Nguyen2023PRXQuantum}、TLSG编码\cite{Pan2025Triangular}）系统化为结构化参考，包括物理动机、小工具构造技巧、正确性论证。

\textbf{逐步推理引导：}引导LLM逐步推理规约设计，从"需要编码什么约束"到"如何用目标哈密顿量结构表达"再到"如何保证基态对应关系"。

\textbf{多模态输入：}对于几何约束（如晶格结构、原子排布），探索将图形表示与文本描述结合的多模态输入方法。

\subsection{ILP验证框架}

申请人团队在三角晶格编码工作中\cite{Pan2025Triangular}首次将小工具验证形式化为整数线性规划(ILP)问题。本项目将扩展这一方法：

\textbf{ILP模型构造：}将规约正确性条件编码为整数线性约束。决策变量对应晶格位置的占用状态（$x_i \in \{0, 1\}$），约束条件编码独立集性质（相邻位点不同时占据）和逻辑线传递条件（信息在小工具间正确传递）。对于给定的小工具构型，ILP可精确判定解保持性和解可提取性。

\textbf{反例驱动改进：}当ILP验证失败时，求解器生成具体的反例构型（即导致正确性违反的具体问题实例），这些反例可直接反馈给LLM指导规约方案的修正。

\subsection{进化式优化}

借鉴AlphaEvolve\cite{AlphaEvolve2025}的成功经验，对已验证规约进行进化优化：

\textbf{进化循环：}初始化（LLM生成种群）$\to$ 评估（ILP验证+代价计算）$\to$ 选择（保留高质量候选）$\to$ 变异（LLM生成变体）$\to$ 迭代。

\textbf{多目标优化：}同时优化多个物理相关指标，包括量子比特开销、质量因子$Q$与能量分离度（针对里德堡编码）、链长度与链断裂率（针对D-Wave嵌入）等。

\textbf{与AlphaEvolve的关键区别：}AlphaEvolve主要应用于数值优化，而本项目针对离散组合结构。关键创新在于将ILP求解器作为精确验证器嵌入进化循环，确保产出规约的可证明正确性。

\subsection{预期成果}

本研究预期在项目周期内完成100+规约规则的发现与验证，覆盖至少15类主流NP难问题（MIS、SAT、MAX-CUT、图着色、集合覆盖、Hamilton回路等）到两类量子优化原生问题（单位圆盘图MIS、QUBO）的映射。每条规约均配备：（i）可执行的映射算法与解提取算法；（ii）ILP验证的数学正确性证明；（iii）精确的代价多项式（量子比特开销关于问题规模的函数）。在方法论层面，建立LLM驱动的测试驱动规约发现方法论，预期发表方法论文2--3篇。同时发布开源规约发现工具包，集成ILP验证。


\section{研究内容三：问题规约编译器设计与实现}

\subsection{系统定位}

本项目的核心产出是一个\textbf{开源问题规约编译器}，为量子优化实验研究者提供自动化映射工具。用户只需指定待求解的组合优化问题实例和目标量子硬件平台（里德堡原子阵列或D-Wave退火机），编译器即可自动完成问题识别、规约路径选择、逐步映射和目标哈密顿量参数生成的全流程。

\subsection{系统架构}

编译器由四个核心模块组成：

\textbf{模块1：规约知识库。}存储研究内容二发现的100+规约规则。每条规则包含规约映射$\phi$和解提取映射$\psi$的可执行实现、代价多项式的精确表达式、正确性验证证据（ILP验证结果或形式化证明），以及适用的硬件后端标注。知识库设计为可增量更新的结构，便于持续积累新发现的规约。

\textbf{模块2：路径规划器。}基于研究内容一的理论实现最优规约路径搜索。给定用户问题类型和目标平台，规划器查询规约图计算所有可行路径，根据问题规模参数选择最优路径，或在不存在全局最优时呈现Pareto前沿供用户决策。

\textbf{模块3：规约执行器。}沿选定路径逐步应用规约规则，负责管理中间问题实例的表示转换，记录每步规约的映射关系以便最终解提取，并支持中间结果的缓存和复用以提高大规模问题的处理效率。

\textbf{模块4：硬件参数生成器。}针对目标平台生成实验所需的物理参数。里德堡后端输出原子阵列坐标和激光脉冲参数（兼容Bloqade等主流实验控制平台）；量子退火后端输出QUBO耦合矩阵$J_{ij}$、局域场$h_i$和嵌入方案（兼容Ocean等主流退火工具）。后端设计遵循统一接口规范，便于适配不同厂商的量子优化平台。

\subsection{关键技术}

\textbf{问题识别：}自动识别用户输入问题的类型和结构参数。支持标准格式（如DIMACS CNF、DIMACS图格式）的直接解析，以及自定义问题的结构化描述。

\textbf{增量更新：}当规约知识库新增规约规则时，支持增量重规划而无需从头计算所有路径，提高实验迭代效率。

\textbf{解的反向提取：}系统不仅完成问题到哈密顿量的正向映射，还支持从量子硬件测量结果反向提取原问题的解。沿规约路径逆向应用解提取映射$\psi$，最终返回原始问题的构型。

\subsection{硬件后端优化}

对于\textbf{里德堡原子阵列}，后端支持多种晶格结构（正方、三角、蜂窝）并自动选择最优编码方案。优化目标是最大化质量因子$Q$和能量分离度（编码基态与首个激发态的能隙），以增强绝热演化中的基态保护。编码过程自动适配硬件参数（原子数上限、阻塞半径$R_b$、激光波长限制），并利用申请人开发的UnitDiskMapping.jl库进行编码优化。

对于\textbf{D-Wave量子退火机}，后端针对Pegasus/Zephyr拓扑结构优化次嵌入方案，以最小化物理量子比特开销和链长度。同时自动调优链强度参数$J_c$，在链完整性和有效能量尺度之间取得最佳平衡，降低链断裂率。

\subsection{预期成果}

本研究预期发布开源问题规约编译器，支持从主流NP难问题到里德堡原子阵列和量子退火机等量子优化平台的端到端自动映射。核心性能指标：千顶点规模问题的映射时间不超过分钟量级，生成方案的量子比特开销和能隙特性不劣于手工设计方案。系统支持自定义规约规则和硬件后端扩展，便于适配新型量子优化平台。


\section{研究内容四：量子优化平台实验验证}

\subsection{验证目标}

规约编译器的价值最终体现在实际量子硬件上的求解效果。本研究将在代表性量子优化平台上开展系统性实验验证，具体平台根据项目进展和硬件可用性选择。编译器设计支持里德堡原子阵列和量子退火机两类后端，实验验证将择其一重点开展。

\subsection{里德堡原子阵列验证方案}

里德堡原子阵列是验证规约编译器的理想平台，其原生问题（单位圆盘图MIS）与编译器的核心输出直接对应。当前代表性平台（如QuEra Aquila）已达到256+量子比特规模，且持续快速升级。申请人在该方向具有深厚积累，曾在哈佛大学Lukin团队从事里德堡量子优化实验研究，并与QuEra Computing保持长期技术合作。

\textbf{验证内容：}首先在真实硬件上验证三角晶格编码\cite{Pan2025Triangular}的性能并与KSG编码\cite{Nguyen2023PRXQuantum}进行系统对比；随后测试LLM发现的新型小工具在实际量子系统中的表现；最后针对同一问题的不同规约路径，比较求解成功率和资源开销，检验路径优化理论的实际效果。目标问题涵盖图着色、MAX-CUT、调度问题等经典NP难问题，规模从数十到数百顶点。

\subsection{量子退火机验证方案}

量子退火机（如D-Wave Advantage系列，5000+量子比特）是验证QUBO规约链的代表性平台，可通过云服务获取访问。

\textbf{验证内容：}重点比较编译器生成的嵌入方案与平台默认嵌入工具在嵌入质量上的差异，评估自动调优的$J_c$参数对链断裂率的改善效果，并在1000+变量的实际问题上测试求解质量。目标问题为QUBO形式的组合优化问题，包括MAX-CUT、调度、投资组合优化等。

\subsection{验证指标}

验证将围绕四个维度展开：\textbf{正确性}——量子硬件的测量结果是否对应原问题的有效解，近似比（approximation ratio）与最优解的差距；\textbf{效率}——相比人工设计的映射方案，编译器生成方案的量子比特开销和求解成功概率是否更优；\textbf{自动化程度}——从问题输入到获得硬件测量结果的端到端时间和人工干预程度；\textbf{可扩展性}——随问题规模增大，所需量子比特数和映射时间的增长趋势。

\subsection{预期成果}

本研究预期在代表性量子优化平台（里德堡原子阵列或量子退火机）上完成至少5类NP难问题、3个以上规模梯度的系统性实验验证。验证编译器生成方案在正确性、资源效率和自动化程度三个维度上达到或超越手工方案。预期发表实验论文2--3篇。


\section{特色与创新点}

本项目的研究具有以下三个方面的特色与创新：

\textbf{创新点一：规约代价模型理论与路径优化框架（理论创新，对应S1）。}与经典编译器优化中使用常数代价模型不同，本项目将规约代价建模为问题规模的多变量多项式，将规约选择形式化为多项式代价图上的路径优化问题，研究代价比较的可判定性和Pareto最优路径的计算。这一理论框架为量子编译优化提供了新的形式化工具。将在规约图规模（100+节点）和路径搜索效率两个维度上，与穷举搜索和贪心策略进行对比验证。

\textbf{创新点二：LLM驱动的测试驱动规约发现方法（方法创新，对应S2）。}与AlphaEvolve\cite{AlphaEvolve2025}等依赖数值评估的LLM辅助算法发现方法不同，本项目将整数线性规划(ILP)求解器作为精确验证器嵌入LLM生成循环，实现从"统计验证"到"数学证明"的正确性保障升级。与纯随机搜索和人工设计基线对比，评估规约发现的效率（发现速度）和质量（代价多项式最优性）。

\textbf{创新点三：面向量子优化平台的端到端规约编译器及开源生态（系统创新，对应S3/S4）。}与当前实验研究者手工设计映射方案的做法不同，本项目构建自动化规约编译器，集成规约知识库、路径规划和硬件参数生成，实现从问题描述到量子硬件执行的端到端自动映射。同时扩展problem-reductions库\cite{ProblemReductions2025}至100+规约规则的开源知识库，每条规约配备代价多项式与正确性证明。以手工设计方案和平台默认工具（D-Wave Ocean SDK）为基线，在标准基准问题集上验证自动化方案的竞争力。


\section{年度研究计划}

本项目为四年期面上项目，研究计划如表\ref{tab:timeline}所示。

\begin{table}[htbp]
	\MSzh\MSen\selectfont
	\centering
	\caption{年度研究计划}
	\label{tab:timeline}
	\begin{tabular}{|c|p{12cm}|}
		\hline
		& \\[-6pt]
		\textbf{年度} & \textbf{主要研究内容与预期成果} \\
		\hline
		& \\[-6pt]
		第一年 &
		\textbf{理论基础与方法原型} \newline
		$\bullet$ 建立规约代价模型的形式化框架，完成代价比较可判定性分析~(S1) \newline
		$\bullet$ 开发LLM驱动规约发现的原型系统并验证方法有效性~(S2) \newline
		$\bullet$ 完成三角晶格编码在里德堡原子阵列平台上的实验验证~(S4) \newline
		$\bullet$ 发现并验证20+规约规则，发表1--2篇论文~(S2) \\
		\hline
		& \\[-6pt]
		第二年 &
		\textbf{编译器核心开发} \newline
		$\bullet$ 实现路径优化算法，支持Pareto前沿的增量计算~(S1) \newline
		$\bullet$ 开发编译器核心模块（知识库、路径规划器、规约执行器）~(S3) \newline
		$\bullet$ 扩展规约知识库至50+规则（含ILP正确性证明）~(S2) \newline
		$\bullet$ 发表2篇论文（代价模型理论+规约发现方法）~(S1, S2) \\
		\hline
		& \\[-6pt]
		第三年 &
		\textbf{系统集成与实验验证} \newline
		$\bullet$ 完成编译器全部模块集成，实现端到端自动映射~(S3) \newline
		$\bullet$ 开发里德堡和量子退火硬件后端并对接实际设备~(S3, S4) \newline
		$\bullet$ 在代表性量子优化平台上完成至少3类问题的实验验证~(S4) \newline
		$\bullet$ 扩展规约知识库至80+规则，发表2--3篇论文~(S2) \\
		\hline
		& \\[-6pt]
		第四年 &
		\textbf{优化完善与开源发布} \newline
		$\bullet$ 根据实验反馈优化路径选择策略与硬件参数生成~(S1, S3) \newline
		$\bullet$ 完成100+规约规则的知识库建设（含完整正确性证明）~(S2) \newline
		$\bullet$ 开源发布完整编译器工具链，撰写用户文档与技术指南~(S3) \newline
		$\bullet$ 发表1--2篇论文（含1篇系统或综述论文）~(S1--S4) \\
		\hline
	\end{tabular}
\end{table}

各年度研究内容相互衔接：第一年建立理论基础并验证方法可行性，第二年开发编译器核心功能，第三年进行系统集成和硬件验证，第四年优化完善并开源发布。项目预期总计发表高水平论文6--10篇，开源发布问题规约编译器及100+规约规则知识库。这一安排确保项目稳步推进，各阶段成果可检验、可积累。

\newpage
